\documentclass[openany]{book}

\usepackage[margin=1in]{geometry}
\usepackage{amsmath,amsfonts,amsthm, amssymb}
\usepackage{yhmath}
\usepackage{mathrsfs}
\usepackage{mathtools}
\usepackage{xcolor}
\usepackage{graphicx}
\usepackage{comment}
\usepackage{tikz-cd}
\usepackage{quiver}
\renewcommand{\familydefault}{ppl}
\newcommand{\tr}{\text{tr}}
\newcommand{\R}{\mathbb{R}}
\newcommand{\E}{\mathbb{E}}
\newcommand{\Z}{\mathbb{Z}}
\newcommand{\C}{\mathbb{C}}
\newcommand{\F}{\mathbb{F}}
\newcommand{\la}{\langle}
\newcommand{\ra}{\rangle}
\newcommand{\colim}{\text{colim}}
\DeclareMathOperator{\im}{im}
\let\oldemptyset\emptyset
\let\emptyset\varnothing
\newcommand{\tor}{\text{Tor}}
\newcommand{\id}{\text{id}}
\newcommand{\ext}{\text{Ext}}
\newcommand{\ptop}{\text{PTop}}
\newcommand{\pt}{\text{pt}}
\newcommand{\ach}{\text{Ach}}
\newcommand{\Q}{\mathbb{Q}}
\newcommand{\gal}{\text{Gal}}


\usepackage{thmtools,thm-restate}

% Fixing mdframed skip below
% See https://tex.stackexchange.com/a/292090/143086
\usepackage[framemethod=TikZ]{mdframed}
\usepackage{xpatch}
\makeatletter
\xpatchcmd{\endmdframed}
	{\aftergroup\endmdf@trivlist\color@endgroup}
	{\endmdf@trivlist\color@endgroup\@doendpe}
	{}{}
\makeatother

\definecolor{huilightpink}{HTML}{fff2fe}
\definecolor{huidarkpink}{HTML}{d955b7}
\declaretheoremstyle[
	mdframed={
		backgroundcolor=huilightpink,
		linecolor=huidarkpink,
		rightline=false,
		topline=false,
		bottomline=false,
		linewidth=2pt,
		innertopmargin=5pt,
		innerbottommargin=8pt,
		innerleftmargin=8pt,
		leftmargin=-2pt,
		skipbelow=2pt,
		nobreak
	},
	headfont=\normalfont\bfseries\color{huidarkpink}
]{huipinkbox}
\declaretheorem[style=huipinkbox,name=Theorem,within=chapter]{thm}
\declaretheorem[style=huipinkbox,name=Theorem,sibling=thm]{theorem}




\begin{comment}
\definecolor{huilightyellow}{HTML}{fff5d6}
\definecolor{huidarkyellow}{HTML}{fcad03}
\declaretheoremstyle[
	mdframed={
		backgroundcolor=huilightyellow,
		linecolor=huidarkyellow,
		rightline=false,
		topline=false,
		bottomline=false,
		linewidth=2pt,
		innertopmargin=5pt,
		innerbottommargin=8pt,
		innerleftmargin=8pt,
		leftmargin=-2pt,
		skipbelow=2pt,
		nobreak
	},
	headfont=\normalfont\bfseries\color{huidarkyellow}
]{huiyellowbox}
\declaretheorem[style=huiyellowbox,name=Proposition,within=chapter]{prop}
\end{comment}



\definecolor{huilightpurple}{HTML}{faf2ff}
\definecolor{huidarkpurple}{HTML}{912ed9}
\declaretheoremstyle[
	mdframed={
		backgroundcolor=huilightpurple,
		linecolor=huidarkpurple,
		rightline=false,
		topline=false,
		bottomline=false,
		linewidth=2pt,
		innertopmargin=5pt,
		innerbottommargin=8pt,
		innerleftmargin=8pt,
		leftmargin=-2pt,
		skipbelow=2pt,
		nobreak
	},
	headfont=\normalfont\bfseries\color{huidarkpurple}
]{huipurplebox}
\declaretheorem[style=huipurplebox,name=Proposition,within=chapter]{prop}



% \definecolor{huilightpurple}{HTML}{faf2ff}
% \definecolor{huidarkpurple}{HTML}{912ed9}
% \declaretheoremstyle[
% 	mdframed={
% 		backgroundcolor=huilightpurple,
% 		linecolor=huidarkpurple,
% 		rightline=false,
% 		topline=false,
% 		bottomline=false,
% 		linewidth=2pt,
% 		innertopmargin=5pt,
% 		innerbottommargin=8pt,
% 		innerleftmargin=8pt,
% 		leftmargin=-2pt,
% 		skipbelow=2pt,
% 		nobreak
% 	},
% 	headfont=\normalfont\bfseries\color{huidarkpurple}
% ]{huipurplebox}
\declaretheorem[style=huipurplebox,name=Lemma,within=chapter]{lem}


\definecolor{lightpink}{HTML}{f0f6fc}
\definecolor{darkpink}{HTML}{2c72b8}
\declaretheoremstyle[
	mdframed={
		backgroundcolor=lightpink,
		linecolor=darkpink,
		rightline=false,
		topline=false,
		bottomline=false,
		linewidth=2pt,
		innertopmargin=5pt,
		innerbottommargin=8pt,
		innerleftmargin=8pt,
		leftmargin=-2pt,
		skipbelow=2pt,
		nobreak
	},
	headfont=\normalfont\bfseries\color{darkpink}
]{pinkbox}
\declaretheorem[style=pinkbox,name=Definition,within=chapter]{defn}


\definecolor{huilightblue}{HTML}{edf9ff}
\definecolor{huidarkblue}{HTML}{4b79db}
\declaretheoremstyle[
	mdframed={
		backgroundcolor=huilightblue,
		linecolor=huidarkblue,
		rightline=false,
		topline=false,
		bottomline=false,
		linewidth=2pt,
		innertopmargin=5pt,
		innerbottommargin=8pt,
		innerleftmargin=8pt,
		leftmargin=-2pt,
		skipbelow=2pt,
		nobreak
	},
	headfont=\normalfont\bfseries\color{huidarkblue}
]{huiblueblox}
\declaretheorem[style=huiblueblox,name=Example,within=chapter]{example}



% \definecolor{huilightblue}{HTML}{edf9ff}
% \definecolor{huidarkblue}{HTML}{4b79db}
% \declaretheoremstyle[
% 	mdframed={
% 		backgroundcolor=huilightblue,
% 		linecolor=huidarkblue,
% 		rightline=false,
% 		topline=false,
% 		bottomline=false,
% 		linewidth=2pt,
% 		innertopmargin=5pt,
% 		innerbottommargin=8pt,
% 		innerleftmargin=8pt,
% 		leftmargin=-2pt,
% 		skipbelow=2pt,
% 		nobreak
% 	},
% 	headfont=\normalfont\bfseries\color{huidarkblue}
% ]{huiblueblox}
% \declaretheorem[style=huiblueblox,name=Example,within=chapter]{example}

% \declaretheoremstyle[
% 	mdframed={
% 		backgroundcolor=huilightblue,
% 		linecolor=huidarkblue,
% 		rightline=false,
% 		topline=false,
% 		bottomline=false,
% 		linewidth=2pt,
% 		innertopmargin=5pt,
% 		innerbottommargin=8pt,
% 		innerleftmargin=8pt,
% 		leftmargin=-2pt,
% 		skipbelow=2pt,
% 		nobreak
% 	},
% 	headfont=\normalfont\bfseries\color{huidarkblue}
% ]{huiblueblox}
\declaretheorem[style=huiblueblox,name=Problem,within=chapter]{prob}



% \declaretheoremstyle[
% 	mdframed={
% 		backgroundcolor=huilightblue,
% 		linecolor=huidarkblue,
% 		rightline=false,
% 		topline=false,
% 		bottomline=false,
% 		linewidth=2pt,
% 		innertopmargin=5pt,
% 		innerbottommargin=8pt,
% 		innerleftmargin=8pt,
% 		leftmargin=-2pt,
% 		skipbelow=2pt,
% 		nobreak
% 	},
% 	headfont=\normalfont\bfseries\color{huidarkblue}
% ]{huiblueblox}
\declaretheorem[style=huiblueblox,name=Exercise,within=chapter]{exer}
\declaretheorem[style=huipinkbox, name=Corollary, within=chapter]{cor}





















\newcommand{\nirwarnsymbol}{%
	\begin{tikzpicture}[baseline=(x.base)]
		\draw[rounded corners=.01em] (-.05em,-1.07em)rectangle(.05em,.78em);
		\draw[fill=white,rounded corners=1.3] (0,.75em)--(.75em,0)--(0,-.75em)--(-.75em,0)--cycle;
		\draw[line width=0.2mm, line cap=round](-.4em,-1.07em)--(.4em,-1.07em);
		\node(x) at (0,0em) {};
		% Thank you https://tex.stackexchange.com/a/262510
		\draw[
			line cap=but,
			line join=round,
			x=.5em,
			line width=0.5mm,
			y=1*(height("Z")-\pgflinewidth)*(1-sin(10)),
			rotate=-10,
			rounded corners=1.5pt,
		](-0.57, 0.57) -- (0.57, 0.57) -- (-0.57, -0.57) -- (0.57, -0.57);
	\end{tikzpicture}%
}

%%%%%%%%%%%%%%%%%%%%%%%%%%%%%%%%%%%%%%%%%%%% MARGINS
\usepackage{marginnote}
% Thank you https://tex.stackexchange.com/a/472882
% Makes marginnotes always appear on the left, apparently
%
\makeatletter
\long\def\@mn@@@marginnote[#1]#2[#3]{%
	\begingroup
		\ifmmode\mn@strut\let\@tempa\mn@vadjust\else
			\if@inlabel\leavevmode\fi
			\ifhmode\mn@strut\let\@tempa\mn@vadjust\else\let\@tempa\mn@vlap\fi
		\fi
		\@tempa{%
			\vbox to\z@{%
				\vss
				\@mn@margintest
				\if@reversemargin\if@tempswa
						\@tempswafalse
					\else
						\@tempswatrue
				\fi\fi

					\llap{%
						\vbox to\z@{\kern\marginnotevadjust\kern #3
							\vbox to\z@{%
								\hsize\marginparwidth
								\linewidth\hsize
								\kern-\parskip
								%\mn@parboxrestore
								\marginfont\raggedleftmarginnote\strut\hspace{\z@}%
								\ignorespaces#1\endgraf
								\vss
							}%
							\vss
						}%
						\if@mn@verbose
							\PackageInfo{marginnote}{xpos seems to be \@mn@currxpos}%
						\fi
						\begingroup
							\ifx\@mn@currxpos\relax\else\ifx\@mn@currpos\@empty\else
									\kern\@mn@currxpos
							\fi\fi
							\ifx\@mn@currpage\relax
								\let\@mn@currpage\@ne
							\fi
							\if@twoside\ifodd\@mn@currpage\relax
									\kern-\oddsidemargin
								\else
									\kern-\evensidemargin
								\fi
							\else
								\kern-\oddsidemargin
							\fi
							\kern-1in
						\endgroup
						\kern\marginparsep
					}%
			}%
		}%
	\endgroup
}
\makeatother
%
% Mostly for todonotes
\renewcommand{\marginpar}{\marginnote}
%%%%%%%%%%%%%%%%%%%%%%%%%%%%%%%%%%%%%%%%%%%% /MARGINS

\definecolor{nirlightred}{RGB}{250, 220, 220}
\definecolor{nirdarkred}{HTML}{f40000}
\declaretheoremstyle[
	mdframed={
		backgroundcolor=nirlightred,
		linecolor=nirdarkred,
		rightline=false,
		topline=false,
		bottomline=false,
		linewidth=2pt,
		innertopmargin=5pt,
		innerbottommargin=8pt,
		innerleftmargin=8pt,
		leftmargin=-2pt,
		skipbelow=2pt,
		nobreak
	},
	headfont=\normalfont\bfseries\color{nirdarkred}
]{nirredbox}

% \makeatletter
% \declaretheorem[
% 	style=nirredbox,
% 	name=Warning,
% 	sibling=thm,
% 	% without \leavevmode, the first item in a list gets misformatted
% 	postheadhook={\leavevmode\marginnote{\nirwarnsymbol}[-3pt]%
% 	\ifthmt@thisistheone% restatable makes alignment weird
% 		\hspace{-2.2pt}%
% 	\fi}
% ]{warn}
% \makeatother

\newcommand{\nirideasymbol}{%
	\begin{tikzpicture}[baseline=(x.base)]
		\draw[rounded corners=.01em] (-.05em,-1.07em)rectangle(.05em,.78em);
		\draw[fill=white,rounded corners=1.3] (0,.75em)--(.75em,0)--(0,-.75em)--(-.75em,0)--cycle;
		\draw[line width=0.2mm, line cap=round](-.4em,-1.07em)--(.4em,-1.07em);
		\node(x) at (0,0em) {};
		\node at (0,0em) {{\textbf{!}}};
	\end{tikzpicture}%
}
\renewcommand{\nirwarnsymbol}{%
	\begin{tikzpicture}[baseline=(x.base)]
		\draw[rounded corners=.01em] (-.05em,-1.07em)rectangle(.05em,.78em);
		\draw[fill=white,rounded corners=1.3] (0,.75em)--(.75em,0)--(0,-.75em)--(-.75em,0)--cycle;
		\draw[line width=0.2mm, line cap=round](-.4em,-1.07em)--(.4em,-1.07em);
		\node(x) at (0,0em) {};
		% Thank you https://tex.stackexchange.com/a/262510
		\draw[
			line cap=but,
			line join=round,
			x=.5em,
			line width=0.5mm,
			y=1*(height("Z")-\pgflinewidth)*(1-sin(10)),
			rotate=-10,
			rounded corners=1.5pt,
		](-0.57, 0.57) -- (0.57, 0.57) -- (-0.57, -0.57) -- (0.57, -0.57);
	\end{tikzpicture}%
}
\makeatletter
\declaretheorem[
	style=nirredbox,
	name=Idea,
	sibling=thm,
	% without \leavevmode, the first item in a list gets misformatted
	postheadhook={\leavevmode\marginnote{\nirideasymbol}[-3pt]%
	\ifthmt@thisistheone% restatable makes alignment weird
		\hspace{-2.2pt}%
	\fi}
]{idea}

\declaretheorem[
	style=nirredbox,
	name=Warning,
	sibling=thm,
	% without \leavevmode, the first item in a list gets misformatted
	postheadhook={\leavevmode\marginnote{\nirwarnsymbol}[-3pt]%
	\ifthmt@thisistheone% restatable makes alignment weird
		\hspace{-2.2pt}%
	\fi}
]{warn}
\makeatother

\title{Calc III Section Notes with Answers
\\ 
\vspace{0.4cm}
\large Fall 2025}




\date{\today}
\author{Hui Sun}


\begin{document}

\maketitle

% \tableofcontents
\newpage


\begin{center}
    \Large Calc III-Week 1 (8/25-29)
\end{center}

\renewcommand\thesection{\arabic{section}}
\noindent
\textbf{Logistics}


\begin{itemize}
    \item TA: Hui.
    \item Email: hsun95@jh.edu.
    \item Office Hours: Tuesday 4-6 PM, Krieger 211; Friday 1-2 PM Zoom.
    \item Biweekly Quizzes: 15 min, 10\%.
    \item Attendance: 5\%. (If you can't make it, email me).
\end{itemize}

\noindent
\textbf{Icebreaking Activity}

\begin{itemize}
    \item In a group of three or four: 
    \begin{enumerate}
        \item Learn each other names, year, pronouns.
        \item Find something in common and different among you and share with the entire class.
        \item Play Buzz if you have time, with prime $7$: say the number if it doens't contain or is not divisible by $7$, say buzz otherwise.
    \end{enumerate}
\end{itemize}


\begin{prob}
    Draw the following vectors in $\R^2$:
    \begin{equation*}
        u=(1,2), \quad v=(3,-2)
    \end{equation*}
    Compute $u+v, u-v$, and draw them in the plane.
\end{prob}
\begin{proof}
    \begin{equation*}
        u+v=(4,0), \quad u-v=(-2, 4)
    \end{equation*}
\end{proof}

\begin{prob}
    Consider the following vectors in $\R^3$:
    \begin{equation*}
        u=(1,2,3), \quad, v=(-2, 1, 4)
    \end{equation*}
    \begin{enumerate}
        \item Compute their norms.
        \item  Two vectors $a,b\in\R^3$ are called \textbf{orthognal} if $a\cdot b=0$. Are $u,v$ orthogonal? If not, find a nonzero vector orthogonal to $u$.
    \end{enumerate}
\end{prob}
\begin{proof}
    \begin{enumerate}
        \item \begin{equation*}
            \|u\|=(u\cdot u)^\frac{1}{2}=\sqrt{14}, \quad \|v\|=\sqrt{21}
        \end{equation*}
        \item We check 
        \begin{equation*}
            u\cdot v=-2+2+12=12\neq 0
        \end{equation*}
        thus not orthogonal. A vector that is orthognal to $u$: $(-3,0,1)$. Note that this vector is \textbf{not} unique! For example, $(-1,-1,1)$ is another such vector.
    \end{enumerate}
\end{proof}

% \begin{prob}
%     \begin{enumerate}
%         \item[(a)] Prove that for \(\vec{x} \in \mathbb{R}^n\) and \(s < t\), \(D_s(\vec{x}) \subset D_t(\vec{x})\).
    
%         \item[(b)] Prove that if \(U\) and \(V\) are neighborhoods of \(\vec{x} \in \mathbb{R}^n\), then so are their union and intersection i.e. show that \(U \cap V\) and \(U \cup V\) are neighborhoods of \(\vec{x} \in \mathbb{R}^n\).
    
%         \item[(c)] Prove that the boundary points of an open interval \((a, b) \subset \mathbb{R}\) are the endpoints \(a\) and \(b\).
%     \end{enumerate}
% \end{prob}
% \begin{proof}
%     \begin{enumerate}
%         \item[(a)] To show $D_s(x)\subset D_t({x})$, it suffices to show that if $y\in D_s(x)$, then $y\in D_t(x)$. By definition,  $D_s(x)=\{z: \|z-x\|<s\}$ and $D_t(x)=\{w: \|w-x\|<t\}$. If $y\in D_s(x)$, then $\|y-x\|<s<t$, hence by definition, $y\in D_t(x)$.
%         \item[(b)] By definition, an neighborhood around $x$ is an open set that contains $x$.
%         \begin{enumerate}
%             \item[(i)] We want to show $U\cup V$ is a neighborhood of $x$, i.e., $U\cup V$ is an open set containing $x$. To check that $U\cup V$ contains $x$, we note $U$ and $V$ both contain $x$, so $x\in U\cup V$ follows.
            
%             It remains to show that $U\cup V$ is an open set. Let $y\in U\cup V$, so we want to show there exists some $r>0$ such that $D_r(y)\subset U\cup V$. Because $y\in U\cup V$, either $y\in U$ or $y\in V$, so there are two cases.
%             \begin{itemize}
%                 \item If $y\in U$, then because $U$ is open, there is $r_1>0$ such that $D_{r_1}(y)\subseteq U$. Thus, $D_{r_1}(y)\subseteq U\cup V$ as well.
%                 \item Otherwise if $y\in V$, then again because $V$ is open, there is $r_2>0$ such that $D_{r_2}(y)\subseteq V$. Thus, $D_{r_2}(y)\subseteq V$ as well.
%             \end{itemize}
%             % Let $y\in U\cup V$ be arbitrary, then either $y\in U$ or $y\in V$. If $y\in U$, then there exists $r_1>0$ such that $D_{r_1}(y)\subset U\subset U\cup V$, and similarly if $y\in V$, then there exists $D_{r_2}(y)\subset V\subset U\cup V$. In either case, we have found some $r_i>0, i=1,2$, such that $D_{r_i}\subset U\cup V$, showing that $U\cup V$ is an open set around $x$ by definition.

%             \item[(ii)] We want to show $U\cap V$ is a neighborhood of $x$. To show $x\in U\cap V$, recall we know $x\in U$ and $x\in V$, so $x\in U\cap V$.
            
%             It remains to show that $U\cap V$ is an open set. This is similar to the previous part. Let $y\in U\cap V$ be arbitrary, and we need to find $r>0$ such that $D_r(y)\subseteq U\cap V$. Because both $y\in U$ and $y\in V$, there are positive real numbers $r_1$ and $r_2$ for which
%             \[D_{r_1}(y)\subseteq U\qquad\text{and}\qquad D_{r_2}(y)\subseteq V.\]
%             Now, let $r=\min\{r_1,r_2\}$, and we claim that $D_r(y)\subseteq U\cap V$, which will completes the proof. Well, following the same proof as in (a), one can show that $r\leq r_1$ implies that $D_r(y)\subseteq D_{r_1}(y)$; similarly, $r\le r_2$ implies that $D_r(y)\subseteq D_{r_2}(y)$. Thus,
%             \[D_r(y)\subseteq D_{r_1}(y)\cap D_{r_2}(y)\subseteq U\cap V,\]
%             as desired.
%             % , then $D_s(x)\subseteq D_t(x)$. Thus $D_r(y)\subseteq D_{r_1}(y)\subset U$ and $D_r(y)\subseteq D_{r_2}(y)\subset V$ because $r\leq r_1$ and $r\leq r_2$. This implies that 
%             % \begin{equation*}
%             %     D_r(y)\subset U\cap V
%             % \end{equation*}
%             % as desired.
%         \end{enumerate}
%         \item[(c)] There are three scenarios we need to consider: (i) elements in $(a,b)$, (ii) endpoints $a$ and $b$, (iii) points outside of $[a,b]$. Recall that a boundary point $x$ of $U$ is defined as a point $x$ such that any neighborhood of $x$ contains one element in the set $U$ and one element outside of the set of $U$.
%         \begin{enumerate}
%             \item[(i)] Let $x\in (a,b)$, so we would like to show $x$ is not a boundary point. We will show that there is an open neighborhood of $x$ contained in $(a,b)$. Let $r=\min\left\{\frac{x-a}{2},\frac{b-x}{2}\right\}$. Consider the neighborhood of $x$ defined by $\left(x-r,x+r\right)$. This set is open and contained in $(a,b)$, so it is the desired neighborhood of $x$.
%             \item[(ii)] Let $x=a$ or $x=b$. We will show that $x$ is a boundary point. Consider first $x=a$:  let $U$ be any open set around $x$, then there exists some $r>0$ such that $D_r(x)\subset U$. By (a) We see that 
%             \begin{equation*}
%                 D_r(x)=(x-r, x+r).
%             \end{equation*}
%             If $x=a$, then 
%             \begin{equation*}
%                 D_r(a)=(a-r, a+r).
%             \end{equation*}
%             We see $D_r(a)$ contains $\frac{a+\min\{r,b\}}{2}$, which is an element in $(a,b)$; $D_r(a)$ also contains $\frac{a-r}{2}$, which is an element not in $(a,b)$. This completes the proof that $a$ is a boundary point. The exact same proof works for $x=b$.
%             \item[(iii)] Let $x\not\in[a,b]$, then $x<a$ or $x>b$, we want to show $x$ is not a boundary point. We will show that there is a neighborhood of $x$ which does not intersect $(a,b)$. We will take $x<a$ as an example: let $r=\frac{a-x}{2}$, then consider the neighborhood 
%             \begin{equation*}
%                 D_r(x)=(x-r, x+r)
%             \end{equation*}
%             One can see that $D_r(x)$ does not contain any points in $(a,b)$ because each $y\in D_r(x)$ has $y<a$. Thus, $x$ is not a boundary point. The exact same argument works for $x>b$.
%         \end{enumerate}
        
        
        
        
%     \end{enumerate}
% \end{proof}

% \newpage


\begin{prob}
    Let $u,v\in\R^3$, suppose that $u,v$ are orthongal, show that 
    \begin{equation*}
        \|u+v\|^2=\|u\|^2+\|v\|^2
    \end{equation*}
    Bonus: is the converse true? (meaning assuming $ \|u+v\|^2=\|u\|^2+\|v\|^2$, is it true that $u\cdot v=0$?)
\end{prob}
\begin{proof}
    We have 
    \begin{align*}
        \|u+v\|^2&=(u+v)\cdot(u+v)\\
        &=u\cdot u+u\cdot v+v\cdot u+v\cdot v\\
        &=\|u\|^2+\|v\|^2
    \end{align*}
    because $u\cdot v=v\cdot u=0$. The converse is also true: we know by definition that
    \begin{equation*}
        \|u+v\|^2=\|u\|^2+\|v\|^2+2u\cdot v
    \end{equation*}
    given the assumption, we also have 
    \begin{equation*}
        \|u+v\|^2=\|u\|^2+\|v\|^2
    \end{equation*}
    Thus equating them we get 
    \begin{equation*}
        \|u\|^2+\|v\|^2+2u\cdot v=\|u\|^2+\|v\|^2\Rightarrow u\cdot v=0
    \end{equation*}
\end{proof}


\noindent
\textbf{Reminders}
\begin{enumerate}
    \item First HW due this Friday.
    \item First Quiz next Tuesday.
\end{enumerate}





\newpage



\begin{center}
    \Large Calc III-Week 2 (9/1-5)
\end{center}



\renewcommand\thesection{\arabic{section}}
\noindent
Topics: (1) cross product, (2) plane in $\R^3$.

\begin{defn}[cross product]
    Let $a=(a_1,a_2,a_3), b=(b_1,b_2,b_3)$ be vectors in $\R^3$, the cross product of $a,b$ is the vector $a\times b$, 
    \begin{equation*}
        a\times b=\begin{bmatrix}
            i&j&k\\
            a_1&a_2&a_3\\
            b_1&b_2&b_3
        \end{bmatrix}
    \end{equation*}
    where $i,j,k$ are the standard vectors in $\R^3$.
\end{defn}

\begin{defn}[Plane in three dimensions]
    A perpendicular vector and a normal vector uniquely define a plane in $\R^3$: given the plane $\mathcal{P}$ passing containing the point $(x_0,y_0,z_0)$ that has a normal vector $(A,B,C)$ is given by the equation:
    \begin{equation*}
        \mathcal{P}: A(x-x_0)+B(y-y_0)+C(z-z_0)=0
    \end{equation*}
\end{defn}

\begin{prop}
    Here are some properties of the cross product:
    \begin{enumerate}
        \item $a\times b$ is perpendicular to vectors $a,b$.
        \item The length of the cross product is the area of the parallelogram:
        \begin{equation*}
            \|a\times b\|=\|a\|\|b\|\sin\theta
        \end{equation*}
        where $\theta$ is the angle between them. (Compare this with the dot product).
        \item $a\times b=-b\times a$, and $a\times (b+c)=a\times b+a\times c$. Moreover, $a\times b=0$ iff $a,b$ are parallel or either $a$ or $b$ are $0$.
        \item (HW) The cross product is \textbf{not} associative! For example, compute 
        \begin{equation*}
            (i\times i)\times j, \quad i\times (i\times j)
        \end{equation*}
    \end{enumerate}
\end{prop}




\begin{prob}
    Let $\vec{u}=(1,2,3), \vec{v}=(0,1,1)$ be vectors in $\R^3$, compute the area of the parallelogram spanned by these two vectors.
\end{prob}
\begin{proof}
    \begin{equation*}
        u\times v=\begin{bmatrix}
            i&j&k\\
            1&2&3\\
            0&1&1
        \end{bmatrix}=-i-j+k=(-1,-1,1)
    \end{equation*}
    Thus the area of the parallelogram is 
    \begin{equation*}
        \|u\times v\|=\sqrt{3}
    \end{equation*}
\end{proof}

\begin{prob}
    Compute the plane containg all three points:
    \begin{equation*}
        (1,0,2), \quad (2, -1, 0), \quad (-1, 2, 3)
    \end{equation*}
\end{prob}
\begin{proof}
    Let $A=(1,0,2), B=(2, -1, 0), C=(-1, 2, 3)$, then consider two vectors in this plane 
    \begin{equation*}
        AB=(1,-1,-2), AC=(-2,2,1)
    \end{equation*}
    Then taking their cross product we find a normal vector to this plane:
    \begin{equation*}
        AB\times AC=\begin{bmatrix}
            i&j&k\\
            1&-1&-2\\
            -2&2&1
        \end{bmatrix}=3i+3j+0k=(3,3,0)
    \end{equation*}
    Thus using the definition above, and point $A$, we know the formula is given by 
    \begin{equation*}
        3(x-1)+3(y)=0
    \end{equation*}
    One can simplify this to 
    \begin{equation*}
        x+y-1=0
    \end{equation*}
\end{proof}

\noindent
\textbf{Reminders}
HW is due Sunday 11:59PM.

% \begin{proof}
%     Find the cross product of $AB=(1,-1,-2), AC=(-2,2,1)$, the 
% \end{proof}


\newpage



\begin{center}
    \Large Calc III-Week 3 (9/8-9/12)
\end{center}

\renewcommand\thesection{\arabic{section}}

\noindent
Topics: (1) graphing multivariable functions, (2) limits and continuity.

\begin{defn}[graph]
    The \textbf{image} of a function $f: \R^n\to\R^m$ is a subset of $\R^m$,
    \begin{equation*}
        \text{Image}(f)=\{f(x)\in\R^m: x\in\R^n\}
    \end{equation*}
    and the \textbf{graph} of $f$ is a subset of $\R^{n+m}$,
    \begin{equation*}
        \text{Graph}(f)=\{(x,f(x)): x\in\R^n\}
    \end{equation*}
\end{defn}



\begin{defn}[limit]
    Let $f: A\subset\R^n\to\R^m$, where $A$ is open, let $x_0$ be in $A$ or be a boundary point of $A$ and $N$ be a neighborhood of a point $b\in\R^m$. Now let $x$ approach $x_0$, $f$ is said to be \textbf{eventually in $N$} if there exists a neighborhood $U$ of $x_0$ such that 
    \begin{equation*}
        \text{ if } x\in U, \text{ then } f(x)\in N
    \end{equation*}
    If $f$ is eventually in $N$ for \textit{any} neighborhood $N$ around $b$, then the \textbf{limit} of $f$ as $x\to x_0$ exists, denoted as 
    \begin{equation*}
        \lim_{x\to x_0}f(x)=b
    \end{equation*}
\end{defn}

\begin{defn}[limit']
    If the limit exists, then $\lim_{x\to x'}f(x)=b$ is when $x=(x_1, x_2, \dots, x_n)\to x'=(x_1',x_2',\dots, x_n')$ from \textbf{all directions}, and $f(x)$ approaches $b=(b_1,\dots, b_m)$.
\end{defn}



\begin{defn}[continuity]
    Let $f:A\subset\R^n\to\R^m$ is said to be \textbf{continuous} at $x_0\in A$ if 
    \begin{equation*}
        \lim_{x\to x_0}f(x)=f(x_0)
    \end{equation*}
    And $f$ is called continuous if $f$ is continuous at every $x_0\in A$.
\end{defn}

\begin{example}
    The limit doesn't need to exist! For example, let 
    \begin{equation*}
        H(x)=\begin{cases}
            1, x\geq 0\\
            -1, x<0
        \end{cases}
    \end{equation*}
    Note the limit doesn't exist at $x=0$.
\end{example}




\begin{prob}
    For the following functions, find their (1) image, (2) graph, (3) draw their graphs.
    \begin{enumerate}
        \item  Let $f: \R\to\R$, and $f(x)=x^2+1$.
        \item Let $g: \R^2\to \R$, and $g(x,y)=x^2+y^2$.
    \end{enumerate}
\end{prob}
\begin{proof}
    \begin{enumerate}
        \item $\text{Image}(f)=\{x^2+1: x\in\R\}$, and $\text{Graph}(f)=\{(x,x^2+1):x\in\R\}$.
        \item $\text{Image}(g)=\{x^2+y^2, (x,y)\in\R^2\}$, and $\text{Graph}(g)=\{(x,y,x^2+y^2):(x,y)\in\R^2\}$.
    \end{enumerate}
\end{proof}

\begin{prob}
    Compute the following limits:
    \begin{enumerate}
        \item \begin{equation*}
            \lim_{(x,y)\to (0,0)}\frac{\sin xy}{y}
        \end{equation*}
       ( Hint: try writing $\frac{\sin xy}{y}=\frac{\sin xy}{xy}\cdot x$, and recall $\lim_{t\to 0}\frac{\sin t}{t}=1$).
        \item \begin{equation*}
            \lim_{(x,y)\to(0,0)}\frac{e^{xy}-1}{y}
        \end{equation*}
        \item \begin{equation*}
            \lim_{(x,y)\to(0,0)}\frac{(x-y)^2}{x^2+y^2}
        \end{equation*}
    \end{enumerate}
\end{prob}
\begin{proof}
    \begin{enumerate}
        \item Following the hint, we see 
        \begin{equation*}
            \lim_{(x,y)\to(0,0)}\frac{\sin xy}{y}=\lim_{(x,y)\to(0,0)}\frac{\sin xy}{xy}x=\lim_{x\to 0}x=0
        \end{equation*}
        \item This one uses the exact same trick: \begin{equation*}
            \lim_{(x,y)\to(0,0)}\frac{e^{xy}-1}{xy}\cdot y=0
        \end{equation*}
        \item First letting $x\to 0$ along $y=0$, we see the limit is $1$; letting $x=y\to 0$, we see the limit is $0$, thus the limit doesn't exist!
    \end{enumerate}
\end{proof}


\begin{prob}
    Compute the limit of the following functions:
    \begin{enumerate}
        \item \begin{equation*}
            \lim_{(x,y)\to(0,0)}\frac{x}{x+y} 
        \end{equation*}
        \item \begin{equation*}
            \lim_{(x,y)\to(0,0)}\frac{xy}{x+y}
        \end{equation*}
        (Hint: try considering $y=x^2-x$ and $y=x$)
        \item \begin{equation*}
            \lim_{(x,y)\to (0,0)}\frac{\sin(xy)}{x+y}
        \end{equation*}
        % \item $f(x,y)=\frac{x}{x+y}$ as $(x,y)\to (0,0)$.
        % \item $g(x,y)=\frac{xy}{x+y}$ as $(x,y)\to (0,0)$.
        % \item $h(x,y)=\frac{\sin(xy)}{x+y}$ as $(x,y)\to (0,0)$. Recall that $\lim_{x\to 0}\frac{\sin(x)}{x}=1$.
    \end{enumerate}
\end{prob}
\begin{proof}
    \begin{enumerate}
        \item First fix $x=0$, let $y\to 0$, then the limit is $0$; now fix $y=0$, let $x\to 0$, the limit is $1$. The limit doesn't exist!
        \item Consider $y=x^2-x$, (as $x\to0, y\to 0$), then 
        \begin{equation*}
            \lim_{(x,y)\to(0,0)}\frac{xy}{x+y}=\lim_{x\to 0}\frac{x^3-x^2}{x^2}=\lim_{x\to 0}x-1=-1
        \end{equation*}
        and consider $y=x$, we see the limit is $0$, thus the limit doesn't exist!
        \begin{warn}
            2 does not follow from 1! A student suggests a proof: $\lim_{(x,y)\to(0,0)}=\frac{x}{x+y}\cdot y$, and by 1, the limit $\frac{x}{x+y}$ doesn't exist, this implies the limit of $\frac{xy}{x+y}$ also doesn't exist. This argument is not correct! Consider the following counterexample: $\lim_{y\to 0}\frac{1}{y}$ doesn't exist, but the limit
            \begin{equation*}
                \lim_{y\to 0}\frac{1}{y}\cdot y=1
            \end{equation*}
            exists! More concretely, if you multiply by any function that doesn't tend to $0$, the argument follows, but it doesn't work when the function tends to $0$! (Sorry I wasn't able to give a concrete counterexample in class other than saying this gives ``bad and untrue vibes''). Thank you (the student) who brought it up, your attempt still remains very very good.
        \end{warn}
        \item We see that 
        \begin{equation*}
            \lim_{(x,y)\to(0,0)}\frac{\sin(xy)}{xy}\frac{xy}{x+y}
        \end{equation*}
        Note that the limit of $\sin(xy)/(xy)=1$, but the second one doesn't exist, thus the limit doesn't exist!
    \end{enumerate}
\end{proof}


\noindent
How to find a a limit $\lim_{x\to x_0}f(x)$:
\begin{itemize}
    \item Step 1: Guess what the limit should be.
    \item Step 2: Try from approaching $x_0$ from different directions.
    \item Step 3: Try to replace terms with expressions you are familiar with.
\end{itemize}



\newpage

\begin{center}
    \Large Calc III-Week 4 (9/15-9/19)
\end{center}

\renewcommand\thesection{\arabic{section}}

\noindent
Topics: (1) Partial derivatives. (2) Definition of total derivatives.

\begin{prob}
    Compute $\frac{\partial f}{\partial x}, \frac{\partial f}{\partial y}$ for the following functions:
    \begin{enumerate}
        \item \begin{equation*}
            x^3y^4-xy^2
        \end{equation*}
        \item \begin{equation*}
            x^2\sin(2y)+3
        \end{equation*}
        \item \begin{equation*}
            \ln\left(\frac{y}{x}\right)+\ln\left(\frac{1}{x+y}\right)-\ln\left(\frac{x}{2}\right)
        \end{equation*}
        You may use the following identities to simply the equation first:
        \begin{equation*}
            \ln\left(\frac{a}{b}\right)=\ln a-\ln b, \quad \ln\left(\frac{1}{a}\right)=-\ln a
        \end{equation*}
    \end{enumerate}
\end{prob}
\begin{proof}
    We have 
    \begin{enumerate}
        \item \begin{equation*}
            \partial x: 3x^2y^4-y^2, \quad \partial y: 4x^3y^3-2xy
        \end{equation*}
        \item \begin{equation*}
            \partial x: 2x\sin(2y), \quad \partial y: 2x^2\cos(2y)
        \end{equation*}
        \item \begin{equation*}
            \partial x: -\frac{2}{x}-\frac{1}{x+y}, \quad \partial y: \frac{1}{y}-\frac{1}{x+y}
        \end{equation*}
    \end{enumerate}
\end{proof}




\begin{defn}[tangent plane]
    Let $f:\R^2\to\R$ be differentiable at $(x_0,y_0)$, then the \textbf{tangent plane} to the graph $f$ in $\R^3$ is the given by 
    \begin{equation*}
        z=f(x_0,y_0)+\frac{\partial f}{\partial x}\bigg\vert_{(x_0,y_0)}(x-x_0)+\frac{\partial f}{\partial y}\bigg\vert_{(x_0,y_0)}(y-y_0)
    \end{equation*}
\end{defn}

\begin{prob}
    Compute the plane tangent to the graph of $f(x,y)=x^2y+2xy-y^2$ at $(1,2)$.
\end{prob}
\begin{proof}
    We have 
    \begin{equation*}
        \frac{\partial f}{\partial x}(1,2)=2xy+2y\vert_{(1,2)}=8, \quad \frac{\partial f}{\partial y}(1,2)=x^2+2x-2y\vert_{(1,2)}=-1
    \end{equation*}
    and $f(1,2)=2$, thus the plane is given by
    \begin{equation*}
        z=2+8(x-1)-(y-2)
    \end{equation*}
    i.e., $z=8x-y-4$.
\end{proof}

\begin{defn}[derivative for two variables]
    Let $f:\R^2\to\R$, then $f$ is said to be \textbf{differentiable} at $(x_0,y_0)$ if $\frac{\partial f}{\partial x}, \frac{\partial f}{\partial y}$ exist at $(x_0,y_0)$ and if 
    \begin{equation*}
        \lim_{(x,y)\to(x_0,y_0)}\frac{f(x,y)-\mathcal{P}(x,y)}{\|(x,y)-(x_0,y_0)\|}=0
    \end{equation*}
    where $\mathcal{P}(x,y)=f(x_0,y_0)+\frac{\partial f}{\partial x}\bigg\vert_{(x_0,y_0)}(x-x_0)+\frac{\partial f}{\partial y}\bigg\vert_{(x_0,y_0)}(y-y_0)$ is the tangent plane to $f$ at $(x_0,y_0)$.
\end{defn}


\begin{defn}[derivative for $n$ variables]
    Let $f:\R^n\to\R$, then the \textbf{gradient} of $f$, denoted as $\nabla f$ is given by 
    \begin{equation*}
        \nabla f=\begin{bmatrix}
            \frac{\partial f}{\partial x_1}, \dots, \frac{\partial f}{\partial x_n}
        \end{bmatrix}
    \end{equation*}
    is a $1\times n$ matrix. And $f$ is said to be \textbf{differentiable} at $x_0\in\R^n$ if 
    \begin{equation*}
        \lim_{x\to x_0}\frac{\|f(x)-f(x_0)-\nabla f(x_0)(x-x_0)\|}{\|x-x_0\|}=0
    \end{equation*}
    and the derivative of $f$ is exactly the gradient $\nabla f$ at $x_0$.
\end{defn}

\begin{defn}[derivative for $m$ outputs]
    Let $f:\R\to\R^m$, where $f(x)=(f_1(x),\dots, f_m(x))$, then let $T$ denote the $n\times 1$ matrix 
    \begin{equation*}
        T=\begin{bmatrix}
            \frac{df_1}{dx}(x_0)\\
            \frac{df_2}{dx}(x_0)\\
            \dots\\
            \frac{df_m}{dx}(x_0)
        \end{bmatrix}
    \end{equation*}
    Then $f$ is said to be \textbf{differentiable} at $x_0$ if 
    \begin{equation*}
        \lim_{x\to x_0}\frac{|f(x)-f(x_0)-T(x-x_0)|}{|x-x_0|}=0
    \end{equation*}
    and the matrix $T$ is the derivative at $x_0$.
\end{defn}
\begin{example}
    Let $f(x)=(x^2, 2x, -x)$, then 
    \begin{equation*}
        T=Df(1)=\begin{bmatrix}
            2x\\
            2\\
            -1
        \end{bmatrix}=
        \begin{bmatrix}
            2\\
            2\\
            -1
        \end{bmatrix}
    \end{equation*}
\end{example}


\begin{defn}[derivative for general functions]
    Let $f: \R^n\to\R^m$, let $T$ be the $m\times n$ matrix with entries $\partial f_i/\partial x_j$ evaluated at $x_0\in \R^n$. Then  $f$ is said to be \textbf{differentiable} at $x_0$ if 
    \begin{equation*}
        \lim_{x\to x_0}\frac{\|f(x)-f(x_0)-T(x-x_0)\|}{\|x-x_0\|}=0
    \end{equation*}
    then $f$ is differentiable at $x_0$, and the matrix $T$ is the derivative at $x_0$. Note that $T$ loosk like 
    \[
T=Df(x_0) , \quad Df=
\begin{bmatrix}
\frac{\partial f_1}{\partial x_1} & \cdots & \frac{\partial f_1}{\partial x_n} \\
\vdots & \ddots & \vdots \\
\frac{\partial f_m}{\partial x_1} & \cdots & \frac{\partial f_m}{\partial x_n}
\end{bmatrix}
\]
\end{defn}
\begin{example}
    Let $f(x,y,z)=(ze^x, -ye^z)$, then 
    \[
Df(x, y, z) = 
\begin{bmatrix}
ze^x & 0 & e^x \\
0 & -e^z & -ye^z
\end{bmatrix}
\]
\end{example}


% \begin{warn}
%     It is not enough to just know the partial derivatives exist at $(x_0,y_0)$. Consider the following example: $f(x,y)=\sqrt{|xy|}$, note that $f$ is not differentiable at $(0,0)$.
% \end{warn}
% \begin{proof}
%     The partial derivatives are 
% \end{proof}







\begin{prob}
    Let \( f : \mathbb{R}^3 \to \mathbb{R} \) be defined by  
\[
f(x, y, z) = x^2 y + y \sin(z) + z e^x.
\]  
Compute the gradient of $f$ at $(1,2,0)$.
\end{prob}
\begin{proof}
    You can compute the partial derivatives 
    \begin{align*}
        \frac{\partial f}{\partial x} &= 2xy + ze^x, \\
        \frac{\partial f}{\partial y} &= x^2 + \sin(z), \\
        \frac{\partial f}{\partial z} &= y \cos(z) + e^x.
        \end{align*}
    thus the gradient is 
    \[
\nabla f = \left( \frac{\partial f}{\partial x}, \frac{\partial f}{\partial y}, \frac{\partial f}{\partial z} \right)=(2xy+ze^x, x^2+\sin z, y\cos z+e^x)
\]
Hence 
\begin{equation*}
    \nabla f(1,2,0)=(4, 1, 2+e)
\end{equation*}

\end{proof}


\newpage

\begin{center}
    \Large Calc III-Week 5 (9/22-9/26)
\end{center}

\renewcommand\thesection{\arabic{section}}

\noindent
Topics: (1) properties of derivatives, (2) directional derivatives, (3) gradient.


% (Section 2.5 on the multivariable chain rule and Section 2.6 on gradients and directional derivatives.

% I would like you to spend about 15-20 minutes going through some chain rule problems e.g. something like problems 7 or 8 from Section 2.5. One problem I really like that reviews multiple different concepts is 11 from Section 2.5 If you could present this problem, I think that would be good for the students
 
% key concept: the matrix of partial derivatives. 
 
% Please use the rest of the time to start discussing Section 2.6, specifically focusing on pg. 135-139. I will try to go through these pages on Monday, but I may have to go quickly. Reviewing all this material with simple examples is what I would recommend. Some problems that could be relevant are 3, 6, and 7 from Section 2.6).


\begin{defn}[path]
    A \textbf{path} $c$ is a map $c: [a,b]\to \R^n$. We can write $c(t)=(c_1(t), \dots, c_n(t))$. If $c$ is differentiable, then we can define the \textbf{velocity} of $c$ at any $t_0\in [a,b]$ as 
    \begin{equation*}
        c'(t_0)=\left(c_1'(t_0), \dots, c_n'(t_0)\right)
    \end{equation*}
    The velocity vector of $c$ at $t_0$ is also a \textbf{tangent} vector to $c$ at $t_0$. The \textbf{speed} of the path $c$ at $t_0$ is the length of the velocity vector $\|c'(t_0)\|$.
\end{defn}

\begin{defn}[tangent line to a path]
    Let $c:[a,b]\to\R^n$ be a path, if $c'(t_0)\neq 0$, then the \textbf{tangent line} at $x_0$ is given by 
    \begin{equation*}
        l(t)=c(t_0)+c'(t_0)(t-t_0)
    \end{equation*}
\end{defn}





\begin{prop}
    Let $f:U\subset\R^n\to\R^m$ be differentiable at $x_0$, then the derivative of $f$ at $x_0$ is an $m\times n$ matrix $Df(x_0)=\left(\frac{\partial f_i}{\partial x_j}\right)_{ij}$. The derivative follows the same properties as derivative for single variable functions:
    \begin{enumerate}
        \item Let $c\in\R$, then 
        \begin{equation*}
            D(cf)(x_0)=cDf(x_0) \tag {multiplication of a matrix by constant $c$}
        \end{equation*}
        \item Let $g: U\subset\R^n\to\R^m$ also be differentiable at $x_0$, then 
        \begin{equation*}
            D(f+g)(x_0)=Df(x_0)+Dg(x_0) \tag{sum of two matrices}
        \end{equation*}
        \item Let $h_1: U\subset\R^n\to\R, h_2: U\subset\R^n\to\R$,then 
        \begin{equation*}
            D(h_1h_2)(x_0)=Dh_1(x_0)h_2(x_0)+h_1(x_0)Dh_2(x_0) \tag{product rule}
        \end{equation*}
        and if $h_2\neq 0$ on $U$.
        \begin{equation*}
            D(h_1/h_2)(x_0)=\frac{Dh_1(x_0)h_2(x_0)-h_1(x_0)Dh_2(x_0)}{h_2^2(x_0)} \tag{quotient rule}
        \end{equation*}
        \item Let $g: U\subset\R^n\to\R^m, f:V\subset\R^m\to\R^p$ such that $g(U)\subset V$, then 
        \begin{equation*}
            D(f\circ g)(x_0)=Df(g(x_0))Dg(x_0) \tag{chain rule}
        \end{equation*}
    \end{enumerate}
\end{prop}

\begin{defn}[directional derivative]
    Let $f:\R^3\to\R$, be differentiable, then the directional directive at $x_0\in\R^3$ in the direction of a \textbf{unit vector} $v$ is given by 
    \begin{equation*}
        \nabla f(x_0)\cdot v=\left[\frac{\partial f}{\partial x_1}(x_0)\right]v_1+\left[\frac{\partial f}{\partial x_2}(x_0)\right]v_2+\left[\frac{\partial f}{\partial x_3}(x_0)\right]v_3
    \end{equation*}
    where $v=(v_1,v_2,v_3)$.
\end{defn}

\begin{prop}
    Suppose that $\nabla f(x_0)\neq 0$, then the direction for which $f$ increases the fastest at $x_0$ is along $\nabla f(x_0)$.
\end{prop}


\begin{prop}
    Let $f:\R^3\to\R$ be differentiable, let $S$ be a level surface of $f$, i.e., $S$ is a surface described by 
    \begin{equation*}
        f(x,y,z)=k
    \end{equation*}
    were $k$ is some constant. Let $(x_0,y_0,z_0)\in S$, then
    \begin{equation*}
        \nabla f(x_0,y_0,z_0) \text{ is normal to the level surface at } (x_0,y_0,z_0)
    \end{equation*}
    This means if $c(t)$ is a path in $S$, and $v(0)=(x_0,y_0,z_0)$, and if $v$ is a tangent vector to $c(t)$ at $t=0$, then 
    \begin{equation*}
        \nabla f(x_0,y_0, z_0)\cdot v=0
    \end{equation*}
    Moreover, if $\nabla f(x_0,y_0,z_0)\neq 0$, the \textbf{tangent plane} of $S$ at $(x_0,y_0,z_0)$ is given by 
    \begin{equation*}
        \nabla f(x_0,y_0,z_0)\cdot (x-x_0, y-y_0, z-z_0)=0
    \end{equation*}
\end{prop}










\begin{prob}
    Consider the curve in $\R$: $c(t)=(2t, t^2,-t)$. Find the speed of the $c$ at $t=2$ and the tangent line at $t=1$.
\end{prob}
\begin{proof}
    The velocity vector of $c$ at $t=2$ is 
    \begin{equation*}
        c'(t)=(2, 2t, -1)
    \end{equation*}
    evaluated at $t=2$ is $c'(2)=(2, 4, -1)$. Thus the speed is the length of the velocity vector 
    \begin{equation*}
        \|c'(2)\|=\left(2^2+4^2+(-1)^2\right)^\frac{1}{2}=\sqrt{21}
    \end{equation*}
    For the tangent line: the tangent vector is 
    \begin{equation*}
        c'(1)=(2,2, -1)
    \end{equation*}
    and $c(1)=(2, 1, -1)$. Thus the tangent line $l$ at $t=1$ is given by 
    \begin{equation*}
        l(t)=(2,1,-1)+t(2,2,-1)=(2+2t, 1+2t, -1-t)
    \end{equation*}
\end{proof}



\begin{prob}[2.5, Q7]
    Let \( f(u, v) = (\tan (u - 1) - e^v, u^2 - v^2) \) and  
   \[
   g(x, y) = (e^{x-y}, x - y).
   \]  
   Calculate \( f \circ g \) and  
   \[
   D(f \circ g)(1, 1).
   \]
\end{prob}
\begin{proof}
    We have 
    \begin{equation*}
        f\circ g(x,y)=(\tan(e^{x-y}-1)-e^{x-y}, e^{2(x-y)}-(x-y)^2)
    \end{equation*}
    and $g(1,1)=(1,0)$, thus using chain rule, we have 
    \begin{equation*}
        D(f\circ g)(1,1)=Df(1,0)Dg(1,1)
    \end{equation*}
    where 
    \begin{equation*}
        Df(u,v)=\begin{bmatrix}
            \sec^2(u-1)& -e^v\\
            2u&-2v
        \end{bmatrix}, \quad Dg(x,y)=\begin{bmatrix}
            e^{x-y}&-e^{x-y}\\
            1&-1
        \end{bmatrix}
    \end{equation*}
    Hence 
    \begin{align*}
        D(f\circ g)(1,1)&=\begin{bmatrix}
            1&-1\\
            2&0
        \end{bmatrix}\begin{bmatrix}
            1&-1\\
            1&-1
        \end{bmatrix}\\
        &=\begin{bmatrix}
            0&0\\
            2&-2
        \end{bmatrix}
    \end{align*}
\end{proof}

\begin{prob}[2.5, Q8]
    Let \( f(u, v, w) = (e^{u-w}, \cos (v + u) + \sin (u + v + w)) \) and \( g(x, y) = (e^x, \cos (y - x), e^{-y}). \)  
   Calculate \( f \circ g \) and \( D(f \circ g)(0, 0)\).
\end{prob}
\begin{proof}
    We have 
    \begin{equation*}
        f\circ g =(e^{e^x-\cos(y-x)}, \cos(e^x+\cos(y-x)), \sin(e^{x}+e^{-y}+\cos(y-x)))
    \end{equation*}
    and $g(0,0)=(1, 1, 1)$. Thus 
    \begin{equation*}
        D(f\circ g)(0,0)=Df(1,1,1)Dg(0,0)
    \end{equation*}
    where 
    \begin{equation*}
        Df(u,v,w)=\begin{bmatrix}
            e^{u-w}&0&-e^{u-w}\\
            -\sin(v+u)+\cos(u+v+w)&-\sin(v+u)+\cos(u+v+w)&\cos(u+v+w)
        \end{bmatrix}
    \end{equation*}
    and 
    \begin{equation*}
        Dg(x,y)=\begin{bmatrix}
            e^x&0\\
            \sin(y-x)&-\sin(y-x)\\
            0&-e^{-y}
        \end{bmatrix}
    \end{equation*}
    Thus 
    \begin{align*}
        D(f\circ g)(0,0)&=\begin{bmatrix}
            1&0&-1\\
            -\sin 2+\cos 3&-\sin 2+\cos 3&\cos 3
        \end{bmatrix}\begin{bmatrix}
            1&0\\
            0&0\\
            0&-1
        \end{bmatrix}\\
        &=\begin{bmatrix}
            1&1\\
            -\sin 2+\cos 3&-\cos 3
        \end{bmatrix}
    \end{align*}
\end{proof}

\begin{prob}[2.5, Q11]
Let \( f(x, y, z) = (3y + 2, x^2 + y^2, x + z^2) \). Let  
\[
c(t) = (\cos(t), \sin(t), t).
\]  

(a) Find the path \( p = f \circ c \) and the velocity vector  
\[
p'(\pi).
\]  

(b) Find \( c(\pi), c'(\pi) \) and \( Df(-1, 0, \pi). \)  

(c) Thinking of \( Df(-1, 0, \pi) \) as a linear map, find  
\[
Df(-1, 0, \pi) \ (c'(\pi)).
\]
\end{prob}
\begin{proof}
    \begin{enumerate}
        \item[(a)] We have 
        \begin{equation*}
            p(t)=(3\sin t+2, 1, \cos t+t^2)
        \end{equation*}
        and 
        \begin{equation*}
            p'(t)=(3\cos t+0, -\sin t+2t)
        \end{equation*}
        thus 
        \begin{equation*}
            p'(\pi)=(-3,0,2\pi)
        \end{equation*}
        \item[(b)] We have $c(\pi)=(-1, 0, \pi)$, and $c'(t)=(-\sin t, \cos t, 1)$, and $c'(\pi)=(0,-1, 1)$. And 
        \begin{equation*}
            Df(x,y,z)=\begin{bmatrix}
                0&3&0\\
                2x&2y&0\\
                1&0&2z
            \end{bmatrix}
        \end{equation*}
        Thus 
        \begin{equation*}
            Df(-1,0,\pi)=\begin{bmatrix}
                0&3&0\\
                -2&0&0\\
                1&0&2\pi
            \end{bmatrix}
        \end{equation*}
        \item[(c)] We have 
        \begin{equation*}
            Df(-1,0,\pi)(c'(\pi))=\begin{bmatrix}
                0&3&0\\
                -2&0&0\\
                1&0&2\pi
            \end{bmatrix}\begin{bmatrix}
                0\\
                -1\\
                1
            \end{bmatrix}=\begin{bmatrix}
                -3\\
                0\\
                2\pi
            \end{bmatrix}
        \end{equation*}
    \end{enumerate}
\end{proof}


\begin{prob}[2.6, Q3]
Compute the directional derivatives of the following functions along unit vectors at the indicated points in directions parallel to the given vector:

% (a) \[ f(x, y) = x^y , \((x_0, y_0) = (e, e)\), \mathbf{d} = 5\mathbf{i} + 12\mathbf{j} \]
(a)
\begin{equation*}
    f(x, y) = x^y, (x_0, y_0) = (e, e), \quad \mathbf{d} = 5\mathbf{i} + 12\mathbf{j}
\end{equation*}

(b)  
\[
f(x, y, z) = e^x + yz, \, (x_0, y_0, z_0) = (1, 1, 1), \quad \mathbf{d} = (1, -1, 1)
\]

(c)  
\[
f(x, y, z) = xyz, \, (x_0, y_0, z_0) = (1, 0, 1), \quad \mathbf{d} = (1, 0, -1)
\]
\end{prob}
\begin{proof}
    We find $\nabla f$ for all these functions and find the directional directive 
    \begin{equation*}
        \nabla f(x_0)\cdot \frac{d}{\|d\|}
    \end{equation*}
    \begin{enumerate}
        \item[(a)] We have 
        \begin{equation*}
            \nabla f(x,y,z)=(yx^{y-1}, x^y\ln x)
        \end{equation*}
        hence 
        \begin{equation*}
            \nabla f(e,e)\cdot \frac{d}{\|d\|}=(e^e,e^e)\cdot\left(\frac{5}{13},\frac{12}{13}\right)=\frac{17}{13}e^e
        \end{equation*}
        \item[(b)] We have 
        \begin{equation*}
            \nabla f(x,y,z)=(e^x, z, y)
        \end{equation*}
        hence 
        \begin{equation*}
            \nabla f(1,1,1)\cdot \frac{1}{\sqrt{3}}(1,-1,1)=\frac{e}{\sqrt{3}}
        \end{equation*}
        \item[(c)] We have 
        \begin{equation*}
            \nabla f(x,y,z)=(yz, xz, xy)
        \end{equation*}
        hence 
        \begin{equation*}
            \nabla f(1,0,1)\cdot \frac{1}{\sqrt{2}}(1,0,-1)=\frac{1}{\sqrt{2}}(0,1,0)\cdot(1,0,-1)=0
        \end{equation*}
    \end{enumerate}
\end{proof}

\begin{prob}[2.6, Q6]
Find a vector which is normal to the curve  
\[
x^3 + xy + y^3 = 11 \text{ at } (1, 2).
\]
\end{prob}
\begin{proof}
    Consider the function $f(x,y)=x^3+xy+y^3$, then the level set of $f(x,y)=11$ coincides with the curve above. Thus it suffices to compute 
    \begin{equation*}
        \nabla f(x,y)=(3x^2+y, x+3y^2)
    \end{equation*}
    and 
    \begin{equation*}
        \nabla f(1,2)=(5, 13)
    \end{equation*}
    is perpendicular to the the level curve.
\end{proof}



\begin{prob}[2.6, Q7]
Find the rate of change of \( f(x, y, z) = xyz \) in the direction normal to the surface  
\[
yx^2 + xy^2 + yz^2 = 3 \text{ at } (1, 1, 1).
\]
\end{prob}





\begin{proof}
    We first find a normal vector to the surface, consider the surface as a level set of the function 
    \begin{equation*}
        g(x,y,z)=yx^2+xy^2+yz^2
    \end{equation*}
    Thus 
    \begin{equation*}
        \nabla g(x,y,z)=(2xy+y^2, x^2+2xy+z^2, 2yz)
    \end{equation*}
    hence 
    \begin{equation*}
        u=\nabla g(1,1,1)=(3, 4, 2)
    \end{equation*}
    is a normal vector to the surface, and we normalize it to get a unit normal vector $n=\frac{u}{\|u\|}=\frac{1}{\sqrt{29}}(3,4,2)$. Now we find the directional derivative of $f(x,y,z)=xyz$ along $(3,4,2)$:
    \begin{equation*}
        \nabla f(x,y,z)=(yz, xz, xy)
    \end{equation*}
    hence 
    \begin{equation*}
        \nabla f(1,1,1)=(1,1,1)
    \end{equation*}
    and the directional derivative is 
    \begin{equation*}
        \nabla f(1,1,1)\cdot n=\frac{9}{\sqrt{29}}
    \end{equation*}
\end{proof}





\noindent 
\textbf{Summary}
\begin{itemize}
    \item If asked to find directional derivative/rate of change along a unit vector $v$ at point $x_0$: find $\nabla f(x_0)\cdot v$.
    \item If asked to find along with direction $f$ increases the fastest: find $\nabla f(x_0)$.
    \item If asked to find a normal vector to a surface at a point $x_0$: construct a function such that the surface is the level set of this function, then find $\nabla f(x_0)$.
\end{itemize}
 




\newpage 


\begin{center}
    \Large Calc III-Week 6 (9/29-10/3)
\end{center}

\renewcommand\thesection{\arabic{section}}

\noindent
Topics: (1) Higher-order derivatives, (2) Taylor expansion.


\begin{prop}
    Let $f(x,y)$ be twice continuously differentiable, then 
    \begin{equation*}
        \frac{\partial^2 f}{\partial x\partial y}=\frac{\partial^2 f}{\partial y\partial x}
    \end{equation*}
\end{prop}


\begin{defn}[First order Taylor expansion]
    Let $f:U\subset\R^n\to\R$ be differentiable at $a\in U$, then 
    \begin{equation*}
        f(x)=f(a)+\sum_{i=1}^n\frac{\partial f}{\partial x_i}(a)(x_i-a_i)+R_1(a,x)
    \end{equation*}
    where 
    \begin{equation*}
        \frac{R_1(a,x)}{\|x-a\|}\to 0 \text{ as } x\to a
    \end{equation*}
\end{defn}

\begin{defn}[Alternative Definition (First-order)]
    Let $f:U\subset\R^n\to\R$ be differentiable at $a\in U$. Then 
    \begin{equation*}
        f(a+h)=f(a)+\sum_{i=1}^nh_i\frac{\partial f}{\partial x_i}(a)+R_1(a,h)
    \end{equation*}
    where $R_1(a,h)/\|h\|\to 0$ as $h\to 0$.
\end{defn}



\begin{defn}[Second order Taylor expansion]\label{taylor}
    Let $f:U\subset\R^n\to\R$ be twice continuously differentiable at $a\in U$, then 
    \begin{equation*}
        f(x)=f(a)+\sum_{i=1}^n\frac{\partial f}{\partial x_i}(a)(x_i-a_i)+\frac{1}{2}\sum_{i,j=1}^n\frac{\partial^2 f}{\partial x_i\partial x_j}(a)(x_i-a_i)(x_j-a_j)+R_2(a,x)
    \end{equation*}
    where 
    \begin{equation*}
        \frac{R_2(a,x)}{\|x-a\|}\to 0 \text{ as } x\to a
    \end{equation*}
\end{defn}


\begin{defn}[Alternative Definition (Second-order)]
    Let $f:U\subset\R^n\to\R$ have continuous partial derivatives of third order. Then we can write 
    \begin{equation*}
        f(a+h)=f(a)+\sum_{i=1}^nh_i\frac{\partial f}{\partial x_i}(a)+\frac{1}{2}\sum_{i,j=1}^nh_ih_j\frac{\partial^2f}{\partial x_i\partial x_j}(a)+R_2(a,h)
    \end{equation*}
    where $R_2(a,h)/\|h\|^2\to 0$ as $h\to 0$.
\end{defn}










\begin{prob}
    Find all the second partial derivatives of $f(x,y)=xy+\ln(x-y)$.

    \noindent
    (This includes $\partial^2f/\partial x^2, \partial^2 f/\partial x\partial y, \partial^2f/\partial y\partial x, \partial^2 f/\partial y^2$).
\end{prob}
\begin{proof}
    For first-order derivatives:
    \begin{equation*}
        \frac{\partial f}{\partial x}=y+\frac{1}{x-y}, \quad \frac{\partial f}{\partial y}=x-\frac{1}{x-y}
    \end{equation*}
    Then 
    \begin{align*}
        \frac{\partial^2f}{\partial x^2}&=-\frac{1}{(x-y)^2}\\
        \frac{\partial^2 f}{\partial x\partial y}&=\frac{\partial^2 f}{\partial y\partial x}=1+\frac{1}{(x-y)^2}\\
        \frac{\partial^2f}{\partial y^2}&=-\frac{1}{(x-y)^2}
    \end{align*}
\end{proof}


\begin{prob}
    Write the second-order Taylor expansion for the following function, 
    \begin{equation*}
        f(x,y)=e^{x+y}
    \end{equation*}
    centered at $(x,y)=(0,0)$.
\end{prob}
\begin{proof}
    We first compute all the partial derivatives:
    \begin{equation*}
        \frac{\partial f}{\partial x}=\frac{\partial f}{\partial y}=\frac{\partial^2f}{\partial x^2}=\frac{\partial^2f}{\partial x\partial y}=\frac{\partial^2f}{\partial y^2}=e^{x+y}
    \end{equation*}
    then 
    \begin{equation*}
        f(x)=1+x+y+\frac{1}{2}x^2+xy+\frac{1}{2}y^2+R_2(x)
    \end{equation*}
    where $R_2(x)/\|x\|\to 0$ as $x\to 0$.
\end{proof}






















\end{document}






















\end{document}